\section{Journal Extension and Evidence Record}
\label{app:extension}
Table~\ref{tab:extension} compares the supplied short manuscript with this revision. It documents the extension of that manuscript; it does not assert a publication status for a workshop version. No new extraction experiment or audit has been added. The complete task-to-deliverable and claim-to-evidence registers are in \artifact{paper/revision/}.

\begin{table}[htbp]
\caption{Workshop-to-journal revision map relative to the supplied short manuscript.}
\label{tab:extension}\footnotesize
\begin{tabularx}{\linewidth}{@{}LLL@{}}\toprule
Existing material & Journal addition & Type of contribution \\\midrule
Five-stage overview & Formal event model, implementation algorithm, stage failures, guarantee table & Method specification \\
Brief time description & Synthetic-time equation, case comparison, timestamp and inclusion manifests & Measurement analysis and reporting correction \\
Category percentages & Activity/context/category distinction and exact transport example & Methodological finding from existing artifacts \\
Coverage caveats & Missing-time proposition with explicit assumptions & New analytical reasoning \\
Short foundry illustration & Transcript timeline, reconstructed event, provenance chain & Evidence-traceability example \\
Condensed audit & Rubric, denominators, ambiguity examples, all judgments & Fuller reporting of existing assessment \\
Numerical results & Portable recalculation, manifests, corrected terminology, six cases & Reproducibility and interpretation \\
Brief limitations & Observability taxonomy, uncertainty path, operational hypotheses, responsible use & Methodological synthesis \\
\bottomrule\end{tabularx}
\end{table}

\section{Reproducibility and Configuration}
\label{app:repro}
The offline command below reads the five frozen artifact directories and the original category map. It writes derived CSVs, a JSON summary, and typeset appendices to a separate output directory. It requires Python 3 and no model credentials or third-party packages.
\begin{lstlisting}
python3 paper/reproduce.py
# Optional explicit locations:
python3 paper/reproduce.py --root /path/to/repository \
  --output-dir /path/to/results
\end{lstlisting}
The output includes \artifact{artifact_manifest.csv} (path, size, SHA-256), \artifact{inclusion_manifest.csv} (all 322 decisions), \artifact{transcript_intervals.csv}, \artifact{structural_metrics.csv}, \artifact{factory_metrics.csv}, \artifact{worker_metrics.csv}, \artifact{worker_comparisons.csv}, \artifact{factory003_transport.csv}, and \artifact{audit_judgments.csv}. The summary records script and prompt hashes. Merged-log consistency is checked separately by reconstructing S5 output in memory and comparing it to every released merged file. Source artifacts remain unchanged.

The original pandas-based assessment and value-analysis scripts accept portable output paths. The original activity-level discovery script retains its miner, noise threshold 0.2, and video case rule. Running it requires the dependencies documented in \artifact{paper/requirements-assessment.txt}. The environment used for this revision is recorded in \artifact{paper/revision/assessment_environment.txt}; it is not claimed to be the original environment. Table~\ref{tab:repro} maps paper results to their calculations.

\begin{table}[htbp]
\caption{Reproduction map. Paths for the original scripts are relative to \artifact{paper/docs/assessment_materials/}.}
\label{tab:repro}\footnotesize
\begin{tabularx}{\linewidth}{@{}LLL@{}}\toprule
Paper result & Calculation and frozen input & Saved evidence \\\midrule
Tables~\ref{main-tab:inventory}, \ref{main-tab:anomalies}, \ref{main-tab:corpus} & \texttt{reproduce.py}; source texts and per-video logs; \texttt{assess.py} for original diagnostics & Artifact, inclusion, interval, and structural manifests \\
Tables~\ref{main-tab:value}, \ref{main-tab:bench} and factory cases & \texttt{reproduce.py} and \texttt{value\_analysis.py}; original process map and retained rows & Factory metrics, worker metrics, comparisons, composition totals \\
Figure~\ref{main-fig:composition} & \texttt{reproduce.py}; exact duration totals and original map & Generated TikZ bars and category percentages \\
Figure~\ref{main-fig:category} and Section~\ref{main-sec:zeros} & Activity-label filter over factory 003 CSVs & 101 rows in \artifact{factory003_transport.csv} \\
Table~\ref{main-tab:audit}, Appendix~\ref{app:audit} & Mechanical transcription of \texttt{audit\_verdicts.md} & 84 judgments in CSV and LaTeX \\
Section~\ref{main-sec:discovery} & Original \texttt{pm\_utility.py}, all per-video rows & \artifact{activity_discovery.json} and execution output \\
Equation~\ref{main-eq:missing} & Algebra under explicitly idealized assumptions & Proof in Section~\ref{main-sec:missingtime}; no empirical bound \\
\bottomrule\end{tabularx}
\end{table}

Regeneration is a separate workflow using scripts 01 to 05 in order. Paths and model identifiers are command-line options in the extraction scripts. S1 requires source video archives and a Gemini credential; S2 and S3 use saved text and Gemini; S4 uses saved text, vocabularies, and an OpenAI credential; S5 is local. Their default model strings are \texttt{gemini-3.1-pro-preview} and \texttt{gpt-5.4}. S4 configures high reasoning effort and a strict JSON response schema; S1 to S3 do not explicitly configure temperature, a seed, or generation-token limits. Transient HTTP status codes 408, 429, 500, 502, 503, and 504 trigger retries; the scripts configure up to five attempts with increasing delay. S1 also has an outer per-video retry loop. These are code settings, not verified execution totals.

The archive does not provide per-request model responses, usage records, exact execution dates, token counts, costs, runtimes, historical dependency versions, or a resolved source revision. Existing outputs may be reused without a content-validity check. Consequently, changing a prompt while retaining existing outputs does not prove that those outputs were generated with the changed prompt. Frozen artifacts, hashes, and the remaining historical queries are recorded so that this limitation is explicit.

\section{Activity-Level Discovery Recalculation}
\label{app:discovery}
\begin{table}[htbp]\centering\caption{Recalculation with the original activity-level discovery settings. Fitness is average trace fitness from token replay; precision is token-based precision.}\label{tab:discovery}
\small\begin{tabular}{@{}lrrr@{}}\toprule Factory & Cases & Fitness & Precision \\ \midrule
001 & 97 & 0.9801 & 0.0691 \\
002 & 54 & 0.9910 & 0.0603 \\
003 & 72 & 0.9874 & 0.0906 \\
004 & 51 & 0.9823 & 0.0666 \\
005 & 22 & 0.9544 & 0.2679 \\
006 & 26 & 0.9720 & 0.0832 \\
\bottomrule\end{tabular}\end{table}
PM4Py version for this recalculation: 2.7.23.11.

The original settings are retained. The saved result records the current PM4Py version. No archived process-level computation or result was found, so no process-level fitness or precision value is reported. Discovery-model images are omitted because the original models were not preserved. The repeated calculation establishes the shown model-log metrics under the recorded current environment, not a recovery of the original runtime.

\section{Complete Operational-Category Map}
\label{app:mapping}
This table reproduces the original map with factory scope made explicit. VA denotes production/completion work under the authors' convention, MH material handling, TR transport, QI inspection, CM cleaning/maintenance, RW rework, and DOC documentation. The sole unmapped label is unused in the archived events.
{\small\begin{longtable}{@{}l p{0.65\linewidth} l@{}}
\toprule Factory & Process label & Category \\ \midrule\endhead
001 & Metal stamping & VA \\
001 & Metal forming & VA \\
001 & Manual mechanical assembly & VA \\
001 & Semi-automated press assembly & VA \\
001 & Structural frame assembly & VA \\
001 & Component bagging & VA \\
001 & Component kitting & VA \\
001 & Material replenishment & MH \\
001 & Finished goods transport & TR \\
001 & Scrap disposal & MH \\
001 & Workstation cleaning & CM \\
001 & Production tracking & DOC \\
001 & Machine maintenance & CM \\
002 & Fabric untangling & MH \\
002 & Fabric edge alignment & VA \\
002 & Pattern matching & VA \\
002 & Fabric folding & VA \\
002 & Fabric pairing & VA \\
002 & Fabric pinning & VA \\
002 & Sub-assembly marking & VA \\
002 & Fabric bundling & VA \\
002 & Garment ironing & VA \\
002 & Garment quality inspection & QI \\
002 & Thread snipping & VA \\
002 & Garment tagging & VA \\
002 & Hanger assembly & VA \\
002 & Garment bagging & VA \\
002 & Cardboard sleeve packaging & VA \\
002 & Cardboard wrapping & VA \\
002 & Package labeling & VA \\
002 & Master bagging & VA \\
002 & Box assembly & VA \\
002 & Box packing & VA \\
002 & Box sealing & VA \\
002 & Cart loading & TR \\
002 & Material transport & TR \\
002 & Floor auditing & DOC \\
003 & Metal casting separation & VA \\
003 & Wax pattern assembly & VA \\
003 & Surface finishing and polishing & VA \\
003 & Defect patching & RW \\
003 & Component sub-assembly & VA \\
003 & Injection molding & VA \\
003 & Mold preparation and cleaning & CM \\
003 & Machine pressing & VA \\
003 & Spot painting & VA \\
003 & Thermal welding & VA \\
003 & Material sorting and sifting & MH \\
003 & Inventory verification & DOC \\
004 & Fabric spreading & VA \\
004 & Bulk fabric cutting & VA \\
004 & Fabric measuring & VA \\
004 & Overlock seaming & VA \\
004 & Continuous fabric chaining & VA \\
004 & Edge binding attachment & VA \\
004 & Piping attachment & VA \\
004 & Curvilinear seam joining & VA \\
004 & Tape feeding & VA \\
004 & Heat press labeling & VA \\
004 & Garment hemming & VA \\
004 & Thread trimming & VA \\
004 & Quality inspection & QI \\
004 & Garment folding & VA \\
004 & Garment ironing & VA \\
004 & Garment bagging & VA \\
004 & Bulk box assembly & VA \\
004 & Pick-and-pack sorting & MH \\
004 & Production data logging & DOC \\
004 & Scrap material management & MH \\
004 & Machine maintenance & CM \\
004 & Material transport & TR \\
005 & motor assembly finishing & unmapped \\
005 & motor labeling and nameplate attachment & VA \\
005 & motor packaging and strapping & VA \\
005 & automated stator winding & VA \\
005 & semi-automatic stator wire insertion & VA \\
005 & stator coil binding & VA \\
005 & stator lead wire finishing and soldering & VA \\
005 & manual lathe machining & VA \\
005 & CNC machine tending & VA \\
005 & drill press machining & VA \\
005 & milling machine operation & VA \\
005 & component deburring and polishing & VA \\
005 & dimensional quality inspection & QI \\
005 & manual hardware pre-assembly & VA \\
005 & batch pneumatic fastening & VA \\
005 & material transport and logistics & TR \\
005 & part sorting and staging & MH \\
005 & finished goods packaging & VA \\
006 & Sand mold creation & VA \\
006 & Pattern cleaning and coating & VA \\
006 & Core and insert placement & VA \\
006 & Mold assembly and closing & VA \\
006 & Mold track advancement & MH \\
006 & Machine cleaning and maintenance & CM \\
006 & Manual production logging & DOC \\
006 & Wooden component scraping & VA \\
006 & Wooden component assembly & VA \\
006 & Machine sanding & VA \\
006 & Powder sifting & VA \\
006 & Chemical dough mixing & VA \\
006 & Dough kneading & VA \\
006 & Machine extrusion & VA \\
006 & Punch press operation & VA \\
006 & Bulk material transport & TR \\
006 & Component staging and arrangement & MH \\
006 & Floor-based part sorting & MH \\
006 & Part mechanical modification & VA \\
006 & Compression molding & VA \\
006 & Raw material mixing & VA \\
006 & Post-processing edge finishing & VA \\
006 & Metal casting separation & VA \\
006 & Workspace clearing & CM \\
\bottomrule\end{longtable}}


\section{Complete Archived Audit Judgments}
\label{app:audit}
The following tables preserve all 84 verdicts and source identifiers. The shorthand \texttt{001/w003/v00009}, for example, identifies factory 001, worker 003, recording 00009. Event verdicts are C (fully supported), P (partially supported), and X (contradicted). Entry verdicts are COV (overlap), EXC (uncovered and judged non-operational), and MISS (uncovered and judged operational). The archive's unsupported truncation explanation is flagged in entry 10. Punctuation is normalized for typesetting; judgments are unchanged.
\subsection{Archived event judgments}
\begin{table*}[!t]\begin{adjustwidth}{-\extralength}{0cm}
\caption{Archived event judgments (part 1 of 2).}\label{tab:audit_event_1}
\footnotesize\setlength{\tabcolsep}{3pt}
\begin{tabularx}{\linewidth}{@{}r p{0.16\linewidth} p{0.20\linewidth} L@{}}
\toprule No. & Source & Interval; verdict & Label and archived note \\ \midrule
1 & \textbf{001}/w003/v00009 & 00:01:38-00:02:12; C & mate components over central jig; assembly over metal block \\
2 & \textbf{001}/w002/v00001 & 00:08:41-00:08:48; C & drop sealed bags into bulk container \\
3 & \textbf{001}/w006/v00000 & 00:07:50-00:08:02; C & receive part from adjacent coworker \\
4 & \textbf{001}/w005/v00006 & 00:17:09-00:17:15; P & actuate machine press; interval covers full load-press-remove cycle \\
5 & \textbf{001}/w005/v00004 & 00:14:00-00:14:55; P & place main component into larger bag; transcript: "repetitive assembly and packaging continues" \\
6 & \textbf{001}/w004/v00000 & 00:06:34-00:06:43; P & drive fastener into metal parts; transcript: "repetitive assembly continues" \\
7 & \textbf{001}/w003/v00006 & 00:04:15-00:06:22; C & insert small pin into base components; split of long composite entry \\
8 & \textbf{002}/w009/v00000 & 00:19:07-00:20:00; C & align printed fabric patterns \\
9 & \textbf{002}/w006/v00002 & 00:05:39-00:06:11; C & enclose garment in pre-printed cardboard folder \\
10 & \textbf{002}/w002/v00007 & 00:10:06-00:10:18; C & attach black plastic hanger hook to bag \\
11 & \textbf{002}/w007/v00002 & 00:08:16-00:08:31; C & wrap white cardboard band around folded shirt \\
12 & \textbf{002}/w005/v00007 & 00:06:09-00:06:24; C & fold t-shirt over removable white template; split of repeated-cycle entry \\
13 & \textbf{002}/w002/v00001 & 00:12:22-00:12:26; P & drop packaged garments into bulk box; transcript: "appearing to drop an item" \\
14 & \textbf{002}/w002/v00001 & 00:11:41-00:11:43; C & drop packaged garments into bulk box \\
15 & \textbf{003}/w003/v00000 & 00:10:36-00:13:22; C & scrape component surface with metal tool \\
16 & \textbf{003}/w005/v00004 & 00:04:13-00:04:29; C & strike metal casting with hammer \\
17 & \textbf{003}/w005/v00004 & 00:15:47-00:15:53; C & scoop loose parts from bulk bin \\
18 & \textbf{003}/w010/v00000 & 00:11:45-00:12:05; C & squeeze binding paste from piping bag \\
19 & \textbf{003}/w001/v00001 & 00:05:45-00:06:12; C & fill metal gap with dark clay \\
20 & \textbf{003}/w010/v00003 & 00:17:12-00:17:29; C & load pre-manufactured insert into mold \\
21 & \textbf{003}/w005/v00000 & 00:03:59-00:07:18; C & strike metal casting with hammer \\
\bottomrule\end{tabularx}\end{adjustwidth}\end{table*}
\begin{table*}[!t]\begin{adjustwidth}{-\extralength}{0cm}
\caption{Archived event judgments (part 2 of 2).}\label{tab:audit_event_2}
\footnotesize\setlength{\tabcolsep}{3pt}
\begin{tabularx}{\linewidth}{@{}r p{0.16\linewidth} p{0.20\linewidth} L@{}}
\toprule No. & Source & Interval; verdict & Label and archived note \\ \midrule
22 & \textbf{004}/w007/v00001 & 00:14:07-00:14:36; C & Iron folded garment with steam iron \\
23 & \textbf{004}/w006/v00004 & 00:16:20-00:16:29; C & Unfold sewn piece to inspect seam tension \\
24 & \textbf{004}/w007/v00001 & 00:07:22-00:07:24; C & Snap large cardboard sheet in half \\
25 & \textbf{004}/w006/v00002 & 00:17:00-00:17:03; C & Snip loose threads with small clippers \\
26 & \textbf{004}/w004/v00005 & 00:14:08-00:14:16; C & Push finished fabric chain to table edge \\
27 & \textbf{004}/w003/v00000 & 00:00:00-00:01:00; C & Hem large fabric piece in short bursts; merges 5 s stand-up interruption \\
28 & \textbf{004}/w005/v00001 & 00:08:44-00:15:18; C & Sew short straight seam; aggregated repetition \\
29 & \textbf{005}/w004/v00002 & 00:04:43-00:04:51; C & brush metal shavings off lathe tool post \\
30 & \textbf{005}/w001/v00000 & 00:06:26-00:06:36; C & position metallic nameplate on motor housing; split of ten-motor loop entry \\
31 & \textbf{005}/w003/v00000 & 00:08:20-00:08:27; C & advance cutting tool via carriage handwheels \\
32 & \textbf{005}/w008/v00001 & 00:19:10-00:19:18; P & arrange parts manually inside deep wire mesh bin; tail of composite transport entry \\
33 & \textbf{005}/w006/v00000 & 00:06:03-00:06:10; C & initiate milling machine cycle via control panel \\
34 & \textbf{005}/w003/v00000 & 00:07:22-00:07:29; C & advance cutting tool via carriage handwheels \\
35 & \textbf{005}/w004/v00000 & 00:02:23-00:02:33; C & apply adhesive label to motor housing \\
36 & \textbf{006}/w003/v00002 & 00:12:28-00:12:38; C & lift empty mold box onto pattern plate; "onto the machine" \\
37 & \textbf{006}/w001/v00004 & 00:05:08-00:06:01; C & sift through mixed small parts on floor; floor bin \\
38 & \textbf{006}/w001/v00002 & 00:08:42-00:08:58; C & guide upper mold half onto lower mold half \\
39 & \textbf{006}/w003/v00003 & 00:11:39-00:11:55; C & pack casting sand into mold frame \\
40 & \textbf{006}/w001/v00002 & 00:11:00-00:11:07; P & guide upper mold half onto lower mold half; transcript vague: "helps with a mold" \\
41 & \textbf{006}/w003/v00003 & 00:04:24-00:04:36; C & level powder surface with wooden scraper \\
42 & \textbf{006}/w003/v00002 & 00:16:15-00:16:30; C & pack casting sand into mold frame \\
\bottomrule\end{tabularx}\end{adjustwidth}\end{table*}
\FloatBarrier
Tables~\ref{tab:audit_event_1} and~\ref{tab:audit_event_2} preserve the event judgments.
\subsection{Archived transcript-entry judgments}
\begin{table*}[!t]\begin{adjustwidth}{-\extralength}{0cm}
\caption{Archived transcript-entry judgments (part 1 of 2).}\label{tab:audit_entry_1}
\footnotesize\setlength{\tabcolsep}{3pt}
\begin{tabularx}{\linewidth}{@{}r p{0.16\linewidth} p{0.20\linewidth} L@{}}
\toprule No. & Source & Interval; verdict & Label and archived note \\ \midrule
1 & \textbf{001}/w008/v00010 & 00:12:00-00:12:51; COV & continues repetitive assembly \\
2 & \textbf{001}/w004/v00007 & 00:05:22-00:05:39; COV & tears open bag of new metal parts \\
3 & \textbf{001}/w002/v00003 & 00:01:07-00:01:12; COV & staples bag shut, sets aside \\
4 & \textbf{001}/w010/v00002 & 00:11:22-00:12:05; COV & aligns hinged-clip rail on plate in jig \\
5 & \textbf{001}/w007/v00004 & 00:15:43-00:16:41; COV & resumes stamping operation \\
6 & \textbf{001}/w008/v00002 & 00:08:46-00:08:49; COV & camera tilt (overlapped by longer event) \\
7 & \textbf{001}/w003/v00005 & 00:09:10-00:09:15; COV & places assembled parts onto cardboard \\
8 & \textbf{002}/w005/v00009 & 00:20:10-00:20:13; COV & places cardboard piece on t-shirt \\
9 & \textbf{002}/w003/v00000 & 00:09:16-00:09:44; COV & fetches and aligns garment stacks \\
10 & \textbf{002}/w010/v00001 & 11:51:00-12:59:00; MISS & aligning stack of pink fabric pieces (archive interprets this as truncation; that cause is unverified) \\
11 & \textbf{002}/w005/v00003 & 00:03:38-00:03:47; COV & irons folded fabric \\
12 & \textbf{002}/w005/v00009 & 00:09:44-00:09:49; COV & irons t-shirt back \\
13 & \textbf{002}/w004/v00006 & 00:13:17-00:13:20; EXC & short walking sequence \\
14 & \textbf{002}/w002/v00009 & 00:11:01-00:12:35; COV & two workers fill sack with garments \\
15 & \textbf{003}/w001/v00002 & 00:18:12-00:18:42; COV & continues scraping task \\
16 & \textbf{003}/w009/v00003 & 00:19:35-00:19:42; COV & empties box onto tray \\
17 & \textbf{003}/w004/v00002 & 00:02:58-00:03:08; COV & applies white paste with brush \\
18 & \textbf{003}/w004/v00007 & 00:19:45-00:19:55; COV & continues filing metal part \\
19 & \textbf{003}/w002/v00000 & 00:18:30-00:18:50; MISS & fetches empty tray, stages tied bundles \\
20 & \textbf{003}/w004/v00002 & 00:06:27-00:06:33; COV & brushes paste onto part \\
21 & \textbf{003}/w002/v00004 & 00:02:48-00:02:51; COV & drops debris into box \\
\bottomrule\end{tabularx}\end{adjustwidth}\end{table*}
\begin{table*}[!t]\begin{adjustwidth}{-\extralength}{0cm}
\caption{Archived transcript-entry judgments (part 2 of 2).}\label{tab:audit_entry_2}
\footnotesize\setlength{\tabcolsep}{3pt}
\begin{tabularx}{\linewidth}{@{}r p{0.16\linewidth} p{0.20\linewidth} L@{}}
\toprule No. & Source & Interval; verdict & Label and archived note \\ \midrule
22 & \textbf{004}/w006/v00004 & 00:17:28-00:17:34; COV & checks seam, folds fabric \\
23 & \textbf{004}/w005/v00003 & 00:00:40-00:00:48; COV & prepares fabric edge for sewing \\
24 & \textbf{004}/w009/v00001 & 00:09:44-00:09:46; COV & drops packaged stack in box \\
25 & \textbf{004}/w006/v00004 & 00:08:12-00:08:29; COV & aligns fabric, resumes sewing \\
26 & \textbf{004}/w003/v00000 & 00:19:50-00:20:00; COV & resumes continuous sewing \\
27 & \textbf{004}/w006/v00004 & 00:11:38-00:11:43; EXC & checks phone briefly \\
28 & \textbf{004}/w006/v00004 & 00:01:24-00:01:40; COV & sews another section \\
29 & \textbf{005}/w002/v00000 & 00:08:04-00:08:33; EXC & brief walk around workbench, returns \\
30 & \textbf{005}/w008/v00000 & 00:01:35-00:01:42; EXC & camera shake while walking \\
31 & \textbf{005}/w003/v00000 & 00:05:32-00:05:40; COV & tightens fixture with wrench \\
32 & \textbf{005}/w008/v00000 & 00:11:36-00:11:45; MISS & standing, inspecting part (borderline operational) \\
33 & \textbf{005}/w011/v00000 & 00:08:56-00:09:22; MISS & presses control-panel buttons, consults sheet \\
34 & \textbf{005}/w008/v00000 & 00:05:32-00:05:54; COV & lifts bin with pallet jack, transports \\
35 & \textbf{005}/w006/v00001 & 00:02:00-00:04:00; COV & workers stack red components into wire basket \\
36 & \textbf{006}/w001/v00002 & 00:08:58-00:09:12; COV & waits with spray gun, sprays mold \\
37 & \textbf{006}/w007/v00000 & 00:03:13-00:03:22; COV & hammers casting apart \\
38 & \textbf{006}/w007/v00000 & 00:15:10-00:15:23; COV & hammers casting, scrapes sand \\
39 & \textbf{006}/w001/v00007 & 00:11:57-00:12:03; COV & blows air onto pattern plate \\
40 & \textbf{006}/w005/v00001 & 00:02:47-00:02:55; COV & carries mold, operates control panel \\
41 & \textbf{006}/w001/v00008 & 00:19:19-00:19:32; COV & flips pattern plate, air-cleans \\
42 & \textbf{006}/w001/v00007 & 00:09:24-00:09:35; COV & two workers move sand mold with lifting clamp \\
\bottomrule\end{tabularx}\end{adjustwidth}\end{table*}
\FloatBarrier
Tables~\ref{tab:audit_entry_1} and~\ref{tab:audit_entry_2} preserve the transcript-entry judgments.


\section{Full Extraction Prompts}
\label{app:prompts}
These are the four released prompt templates, with punctuation normalized for display. The source files and their hashes remain in the resource. They specify requested model behavior, not guaranteed output properties. Placeholder braces are filled by the scripts; additional context is constructed as described in Table~\ref{main-tab:stages}.
\subsection{S1: Video transcription}
\lstinputlisting[basicstyle=\scriptsize\ttfamily]{generated/annotation_prompt.txt}
\subsection{S2: Worker summary}
\lstinputlisting[basicstyle=\scriptsize\ttfamily]{generated/worker_process_summary_prompt.txt}
\subsection{S3: Factory report and vocabularies}
\lstinputlisting[basicstyle=\scriptsize\ttfamily]{generated/factory_process_mining_prompt.txt}
\subsection{S4: Event extraction}
\lstinputlisting[basicstyle=\scriptsize\ttfamily]{generated/annotation_to_event_log_prompt.txt}
